\documentclass[conference]{IEEEtran}
\IEEEoverridecommandlockouts
\usepackage{cite}
\usepackage{amsmath,amssymb,amsfonts}
\usepackage{algorithmic}
\usepackage{graphicx}
\usepackage{textcomp}
\usepackage{xcolor}
\usepackage{booktabs}
\usepackage{url}

\def\BibTeX{{\rm B\kern-.05em{\sc i\kern-.025em b}\kern-.08em
    T\kern-.1667em\lower.7ex\hbox{E}\kern-.125emX}}

\begin{document}

\title{Design, Verification, and Performance Evaluation of a Hybrid Encrypted Messaging Protocol}

\author{\IEEEauthorblockN{1\textsuperscript{st} Your Name Here}
\IEEEauthorblockA{\textit{Department of Computer Science} \\
\textit{Your University Name}\\
City, Country \\
email address}
}

\maketitle

\begin{abstract}
End-to-end encrypted (E2EE) messaging systems require cryptographic protocols that balance rigorous security with computational efficiency. In this paper, we present the design and evaluation of a secure messaging application that utilizes Hybrid Encryption. We implement Elliptic Curve Diffie-Hellman (ECDH) for key encapsulation alongside AES-GCM for symmetric data encryption, enforcing Forward Secrecy through automated key rotation. To validate the security of the protocol against Man-in-the-Middle (MITM) and replay attacks, we provide formal verification using the Scyther tool. Finally, we present an empirical performance benchmark comparing Elliptic Curve Cryptography against legacy algorithms such as RSA and ElGamal across varying data payloads, demonstrating the overwhelming efficiency of ECC for modern scalable systems.
\end{abstract}

\begin{IEEEkeywords}
Cryptography, Elliptic Curve, RSA, ElGamal, Hybrid Encryption, Scyther, Forward Secrecy
\end{IEEEkeywords}

\section{Introduction}
Secure communication is the backbone of modern digital infrastructure. As data privacy concerns grow, platforms have universally adopted end-to-end encryption. However, developers must choose between various cryptographic primitives such as RSA, ElGamal, and Elliptic Curve Cryptography (ECC). 

In this paper, we explore the implementation of a custom E2EE chat application. We verify its security properties mathematically and empirically demonstrate why modern systems favor ECC over RSA and ElGamal through rigorous benchmarking on multi-sized data sets.

\section{System Architecture}
Our architecture utilizes a Hybrid Encryption model to ensure both speed and security.

\subsection{Key Exchange and Encapsulation}
Instead of encrypting bulk data using computationally heavy asymmetric algorithms, we utilize asymmetric cryptography strictly for Key Encapsulation. Our system generates ephemeral key pairs using SECP256R1 and derives a 256-bit symmetric session key via HMAC-based Extract-and-Expand Key Derivation Function (HKDF).

\subsection{Data Encryption}
Bulk data is encrypted using Advanced Encryption Standard in Galois/Counter Mode (AES-GCM). AES-GCM provides both confidentiality and data authenticity, ensuring that ciphertext cannot be maliciously modified in transit.

\subsection{Forward Secrecy}
To mitigate the risk of long-term key compromise, our system implements Perfect Forward Secrecy (PFS). The AES session key is strictly rotated every 5 messages. Old keys are immediately destroyed, permanently sealing past conversations from decryption even if the current key is stolen.

\section{Formal Verification with Scyther}
To ensure the cryptographic protocol is immune to logical flaws, we modeled the protocol using the Security Protocol Description Language (SPDL). We evaluated the model using the Scyther verification tool.

The formal verification proved that our protocol maintains:
\begin{itemize}
    \item \textbf{Secrecy:} The AES session key remains known only to Alice and Bob.
    \item \textbf{Authentication:} Both endpoints are guaranteed non-repudiation and protection against MITM attacks.
\end{itemize}

\section{Performance Benchmarking}
To justify the use of ECC over legacy algorithms, we benchmarked our ECC implementation against custom RSA and ElGamal Key Encapsulation mechanisms. 

\subsection{Methodology}
The algorithms were evaluated on their Key Generation time, Key Encapsulation time, and Data Encryption time. Tests were run on standardized file sizes: 1-Page text (~6 KB), 10-Page text (~33 KB), 50-Page text (~151 KB), and a Multimedia Image (~1 MB). 

\subsection{Results}
Our benchmarking revealed that the symmetric encryption time (AES) scales linearly with the file size, peaking at roughly ~2 MB of memory usage for a 1 MB image. The asymmetric operations remained constant regardless of file size.

As shown in our empirical data, ECC vastly outperformed the legacy alternatives. ECC Key Generation required $<0.001$ seconds, whereas RSA required approximately $0.03$ seconds. For Key Exchange, ElGamal proved highly inefficient requiring nearly $0.04$ seconds due to heavy modular exponentiation. Furthermore, ECC required a public key size of only 356 bytes, compared to 451 bytes for RSA, solidifying ECC as the optimal choice for bandwidth-constrained applications.

\section{Conclusion}
This paper presented the successful implementation and verification of a secure chat protocol. Through formal verification in Scyther, we proved its mathematical security. Through empirical benchmarking against RSA and ElGamal, we demonstrated that ECC provides an unmatched balance of lightweight key sizes and rapid computational speeds, making it the definitive standard for modern encrypted messaging.

\begin{thebibliography}{00}
\bibitem{b1} C. Shannon, ``Communication Theory of Secrecy Systems,'' Bell System Technical Journal, 1949.
\bibitem{b2} W. Diffie and M. Hellman, ``New directions in cryptography,'' IEEE Transactions on Information Theory, 1976.
\bibitem{b3} Cremers, C. J. F. ``The Scyther Tool: Verification, falsification, and analysis of security protocols.'' International Conference on Computer Aided Verification, 2008.
\end{thebibliography}

\end{document}
